\section{任意の滑らかなJacobianブロックの代数的中心三次元}
\label{sec:pg-universal-determinant}
本節では、$\hbar$進の形式的中心元、
相対formalityおよびDolgushevのunimodularity定理を
\emph{使わずに}、任意の三変数Jacobianブロックの中心三次元を
代数的に構成する。
重要なのは、三変数制限
$g=f(x,y,z,0)$の滑らかさではなく、
\emph{全変数に関する三次超曲面$f=0$の滑らかさ}である。
$g$が節点型・尖点型・可約、あるいは恒等的に零の場合も含む。
証明は（i）古典的Casimir核の一定次元性、
（ii）Hesse型で既知の明示中心元、
（iii）普遍交換子行列の小行列式の恒等的消滅、
（iv）Cramer公式からなる。
これはEtingof--Ginzburg\cite[\S8.3]{EG}が
potentialの中心元を特殊化する際に用いた
rank一定性・射影的閉性の考え方と近縁だが、
ここでは固定次数の行列式によって
$\hbar=0$の規格化も直接制御する。

$R=\CC[u_1,\ldots,u_m]$、
$S=R[x,y,z]$、$f\in S_3$とし、
\begin{equation}
 \begin{aligned}
 A_{f,\hbar}
 =R[\hbar]\langle x,y,z\rangle/
 \bigl(&[y,z]-\hbar\sym_R(f_x),\\
       &[z,x]-\hbar\sym_R(f_y),\
       [x,y]-\hbar\sym_R(f_z)\bigr)
 \end{aligned}
 \label{eq:pg-det-quantum}
\end{equation}
とおく。$\sym_R$はactive変数のWeyl対称化である。
$f_{\mathrm{cent}}=f(0,0,0,u)$と書く。

\begin{lemma}[滑らかさとactive偏微分の最大公約因子]
\label{lem:pg-det-gradient-coprime}
$V(f)\subset\PP^{m+2}$が滑らかであるとき、
\begin{equation}
 \gcd(f_x,f_y,f_z)=1
 \quad\text{in }\CC[x,y,z,u_1,\ldots,u_m].
 \label{eq:pg-det-gradient-coprime}
\end{equation}
特に$f$には少なくとも一つactive変数を含む三次単項式が現れる。
\end{lemma}
\begin{proof}
$f$がactive変数に依存しなければ、
$\PP(x,y,z)$上で$f$と全偏微分が零となり、
滑らかさに反する。
$f_x,f_y,f_z$が正次数の共通因子を持つと仮定し、
その既約斉次因子を$h$とする。
$h=0$は$\PP^{m+2}$の次元$m+1$の射影超曲面である。
そこへ残る$m$個の斉次方程式
$f_{u_1}=\cdots=f_{u_m}=0$を課しても、
Krullの高さ定理によりその共通零点の
アフィン錐の次元は少なくとも
\begin{equation}
 (m+3)-(m+1)=2.
 \label{eq:pg-det-gradient-dimension}
\end{equation}
従って零ベクトル以外の共通零点が存在する。
そこでは全ての偏微分が零であり、
Euler等式から$f=0$でもある。
これは$V(f)$の滑らかさに矛盾する。
ゆえに共通非定数因子はない。
\end{proof}

\begin{lemma}[全三次式の古典的Casimir核]
\label{lem:pg-det-classical-kernel}
$V(f)\subset\PP^{m+2}$が滑らかなら、
三変数Jacobian Poisson括弧に対し
\begin{equation}
 \boxed{\quad
 \ker\left(S_3\xrightarrow{p\mapsto
 (\{p,x\}_f,\{p,y\}_f,\{p,z\}_f)}
 S_4^{\oplus3}\right)
 =\CC f\oplus R_3.
 \quad}
 \label{eq:pg-det-classical-kernel}
\end{equation}
$m=0$では$R_3=0$と解釈する。
$f(x,y,z,0)$は特異でも零でもよい。
\end{lemma}
\begin{proof}
$S$はUFDであり、補題
\ref{lem:pg-det-gradient-coprime}により
$(f_x,f_y,f_z)$の最大公約因子は1である。
$p\in S_3$がCasimirなら
\begin{equation}
 d_xp\wedge d_xf=0,
 \qquad
 (p_x,p_y,p_z)=a(f_x,f_y,f_z)
 \quad\text{in }\operatorname{Frac}(S)^3
 \label{eq:pg-det-casimir-proportional}
\end{equation}
となる。
既約因子の付値を比較すると、$p_x,p_y,p_z\in S$かつ
$\gcd(f_x,f_y,f_z)=1$より$a\in S$。
両辺は斉次次数2なので、$d_xp\ne0$なら$a$は
次数零の多項式、すなわち定数$c\in\CC$である。
$d_xp=0$なら$p\in R_3$である。
従っていずれの場合も$p-cf\in R_3$。
逆に$f$と$R$はJacobian括弧のCasimirなので
等号を得る。
$f$はactive変数に依存するから直和である。
\end{proof}

\subsection{普遍交換子行列の階数が跳ばないこと}
$\mathcal T_m\subset S_3$を、
$V(f)\subset\PP^{m+2}$が滑らかな三次式の
非空Zariski開集合とする。
$\mathcal T_m$はアフィン空間の開部分なので既約である。
$\mathcal T_m^\circ$を$f(x,y,z,0)$が滑らかな
平面三次を定める非空開部分とする。
$\mathcal T_m^\circ$も稠密である。

命題\ref{prop:pg-PBW}のrank条件を係数変数も含めて考える。
$S_{g,\hbar}$を定める$3\times3$の線形行列
$M_{g,\hbar}$は係数について多項式的であり、
$\hbar=0$では全非零$x\in\PP^2$でrank2である。
従ってrank1以下のincidenceを射影$\PP^2$から
パラメータ空間へproperに押し出すと閉集合が得られ、
その補集合は
\begin{equation}
 \mathcal U_m\subset\mathcal T_m\times\mathbb A^1_\hbar,
 \qquad
 \mathcal T_m\times\{0\}\subset\mathcal U_m
 \label{eq:pg-det-universal-open}
\end{equation}
というZariski開集合となる。
各ファイバーで三次元AS-regular性と
Berger--TailleferのPBW定理
\cite{PMoriSmith,PBergerTaillefer}が成立する。
固定内部次数$d$の商は有限表示であり、
全ファイバー次元が
$\binom{m+d+2}{d}$に等しい。
ここから各次数部分の局所自由性が従うことを明示する。
底のアフィン開部分上でこの商の有限表示行列を取り、
その行数を$q$とすると、全ての閉点で行列の階数は
$r=q-\binom{m+d+2}{d}$である。
底は被約な有限型$\CC$-スキームなので、
全ての$(r+1)$次小行列式は恒等的に零である。
各点の近傍で非零の$r$次小行列式を可逆にし、
行と列の基本変形を行うと行列は
$\operatorname{diag}(I_r,0)$の形となる。
従ってその余核、すなわち次数$d$部分は局所自由である。
$d=3,4$のWeyl対称化基底は$\hbar=0$で
通常の単項式基底なので、さらに
$\mathcal U_m$を零切断の近傍へ縮めれば、
その基底が全てのファイバーで有効となる。

$N_d=\binom{m+d+2}{d}$、
$d_m=\binom{m+2}{3}$、
$r_m=N_3-d_m-1$とおく。
Weyl座標表示で次の有限行列を得る：
\begin{equation}
 \mathcal D_{f,\hbar}:S_3\longrightarrow S_4^{\oplus3},
 \qquad
 p\longmapsto
 \frac1\hbar\left(
 [\sym_R(p),x],[\sym_R(p),y],[\sym_R(p),z]\right).
 \label{eq:pg-det-commutator-matrix}
\end{equation}
$\hbar=0$で交換子が消え、商が$\hbar$-flatなので
右辺の$1/\hbar$は正則であり、
$\mathcal D_{f,0}$は古典Jacobian括弧を与える。

\begin{lemma}[量子交換子行列の普遍階数恒等式]
\label{lem:pg-det-universal-rank}
$\mathcal U_m$の上で
$\mathcal D_{f,\hbar}$の全ての$(r_m+1)$次小行列式は
恒等的に零である。
特に任意の滑らかな$f$について、
$\hbar=0$を含む開集合上で
\begin{equation}
 \operatorname{rank}\mathcal D_{f,\hbar}\le r_m.
 \label{eq:pg-det-universal-rank}
\end{equation}
\end{lemma}
\begin{proof}
まず任意の閉点$f\in\mathcal T_m^\circ$を固定する。
系\ref{cor:pg-smooth-ternary-explicit}が
その$\hbar=0$近傍で$f$を持ち上げる
\emph{代数的な厳密中心三次元}を与える。
その中心元と$d_m$個の純中心三次単項式は
法$\hbar$で線形独立なので、
必要ならこの近傍を縮めれば、
$\mathcal D_{f,\hbar}$の核の次元は
少なくとも$d_m+1$である。
従って全$(r_m+1)$次小行列式は
既約な開曲線$\mathcal U_m\cap(\{f\}\times\mathbb A^1)$の
非空開部分で零であり、その全体で零となる。
これを全ての$f\in\mathcal T_m^\circ$に適用すると、
小行列式は$\mathcal U_m$の稠密開部分
$\mathcal U_m\cap(\mathcal T_m^\circ\times\mathbb A^1)$で零となる。
この議論にはHesse座標を$f$について代数的に同時選択する必要はない。
これらの小行列式は
$\mathcal U_m$上の\emph{正則関数}であり、
$\mathcal U_m$は既約かつ被約である。
したがって小行列式は$\mathcal U_m$上恒等的に零。
これは特異な三変数制限への
\emph{係数についての恒等式の特殊化}であり、
稠密開上の中心元をそのまま特殊化して
分母が消えないと仮定する議論ではない。
\end{proof}

\begin{theorem}[全滑らかな三次ブロックに対する普遍Cramer中心元]
\label{thm:pg-universal-determinant}
$V(f)\subset\PP^{m+2}$が滑らかな任意の$f\in S_3$に対し、
三変数制限$f(x,y,z,0)$が節点型、尖点型、可約、
または恒等的零であっても、
$0$を含むZariski開集合$U_f\subset\mathbb A^1_\hbar$上に
\begin{equation}
 \boxed{\quad
 \Omega_{f,\hbar}\in
 Z\!\left(A_{f,\hbar}|_{U_f}\right)_3,\qquad
 \Omega_{f,0}=f
 \quad}
 \label{eq:pg-det-all-centers}
\end{equation}
が存在する。
しかも$\Omega_{f,\hbar}$は
\emph{有限個の非可換三次単項式の線形結合}であり、
その係数は$\hbar^2$の有理関数で、
全ての分母は$\hbar=0$で非零となる。
以下のCramer公式が、与えられた$f$から中心元を
有限回の四則演算と行列式計算で構成するアルゴリズムである。

$f$のactive変数を含む三次単項式のうち、
係数$a\ne0$を持つ一つを$w_0$とする。
純中心三次単項式と$w_0$を除いた
$r_m$個の三次単項式を
$w_1,\ldots,w_{r_m}$とする。
$\mathcal D_{f,0}|_{\langle w_1,\ldots,w_{r_m}\rangle}$
は単射なので、$r_m$個の行を選んで
\begin{equation}
 M_{ij}(\hbar)=\ell_i(\mathcal D_{f,\hbar}(w_j)),
 \qquad
 b_i(\hbar)=\ell_i(\mathcal D_{f,\hbar}(w_0)),
 \qquad
 \det M(0)\ne0
 \label{eq:pg-det-minor}
\end{equation}
となるようにする。
$M_j$を$M$の第$j$列を$b$に換えた行列とすれば、
\begin{equation}
 \boxed{\quad
 \Omega_{f,\hbar}
 =f_{\mathrm{cent}}+
 a\left(
 \sym_R(w_0)-\sum_{j=1}^{r_m}
 \frac{\det M_j(\hbar)}{\det M(\hbar)}\sym_R(w_j)
 \right).
 \quad}
 \label{eq:pg-det-universal-center}
\end{equation}
が所要の中心元である。
\end{theorem}
\begin{proof}
補題\ref{lem:pg-det-classical-kernel}によって
$\operatorname{rank}\mathcal D_{f,0}=r_m$。
また、$f$の$w_0$係数は非零なので、
古典Casimir核$\CC f\oplus R_3$と
$W=\langle w_1,\ldots,w_{r_m}\rangle$の交叉は零。
従って$\mathcal D_{f,0}|_W$は単射であり、
指定した$r_m\times r_m$小行列$M$が選べる。
$\det M(0)\ne0$なので、$\hbar=0$を含む開集合へ
縮めれば$M$は可逆となる。

一方、補題\ref{lem:pg-det-universal-rank}により
$\mathcal D_{f,\hbar}$のrankは高々$r_m$。
$M$が可逆なので少なくとも$r_m$でもあり、
従って$M$を構成する$r_m$行は
交換子行列の\emph{全ての行}の行空間を張る。
Cramer公式は
\begin{equation}
 \mathcal D_{f,\hbar}
 \left(w_0-\sum_j\frac{\det M_j(\hbar)}{\det M(\hbar)}w_j\right)=0
 \label{eq:pg-det-cramer-kernel}
\end{equation}
を満たし、これは三つの生成元$x,y,z$との
\emph{全交換子が厳密に零}であることを意味する。
$u_\alpha$は定義から中心なので
$\Omega_{f,\hbar}$は代数全体の中心元である。

$\hbar=0$では括弧の核は$\CC f\oplus R_3$であり、
括弧内の元は$w_0$係数が1、
純中心三次係数が0である。
したがってその特殊化は
$(f-f_{\mathrm{cent}})/a$である。
これが$\Omega_{f,0}=f$を示す。

最後に$\hbar\mapsto-\hbar$と語の反転による
反自己同型（補題\ref{lem:pg-center-parity}）を考える。
Weyl基底の各元は不変であり、
$[\sym_R(p),x_i]/\hbar$のWeyl係数は
$\hbar\mapsto-\hbar$に関して不変。
従って$\mathcal D_{f,\hbar}$、$M,b$と
Cramer公式の係数は$\CC(\hbar^2)$に属する。
全ての操作は次数三・四の有限自由加群で完結し、
形式冪級数も無限次数の中心補正も必要ない。
\end{proof}

\subsection{中心元の規格化と係数空間上の構成}
\label{subsec:pg-det-center-family}
普遍交換子行列の階数がちょうど$r_m$となる開部分を
$\mathcal V_m\subset\mathcal U_m$とする。
補題\ref{lem:pg-det-classical-kernel}により、
$\mathcal V_m$は零切断$\mathcal T_m\times\{0\}$全体を含む。
この上で
\begin{equation}
 \mathcal E_d=S_d\otimes_\CC\OO_{\mathcal V_m},\qquad
 \mathcal C_m=R_3\otimes_\CC\OO_{\mathcal V_m},\qquad
 \rho:\mathcal E_3\longrightarrow\mathcal C_m,
 \quad p\longmapsto p(0,0,0,u)
 \label{eq:pg-det-degree-three-bundles}
\end{equation}
とおき、$\mathcal E_3^{\mathrm{act}}=\ker\rho$とする。
従って$\mathcal E_3=\mathcal C_m\oplus\mathcal E_3^{\mathrm{act}}$である。
$\mathcal A$をこの上の普遍量子代数とし、
以下、Weyl同型によって$\mathcal E_3$と$\mathcal A_3$を同一視する。

\begin{lemma}[割った交換子の核と底変更]
\label{lem:pg-det-divided-kernel}
\begin{equation}
 \mathcal K=\ker\mathcal D,\qquad
 \mathcal L=\ker(\mathcal D|_{\mathcal E_3^{\mathrm{act}}})
 \label{eq:pg-det-kernel-line}
\end{equation}
とおくと、$\mathcal K$は階数$d_m+1$、$\mathcal L$は階数一の
局所自由加群であり、
\begin{equation}
 \mathcal K=\mathcal C_m\oplus\mathcal L,
 \qquad \mathcal K/\mathcal C_m\simeq\mathcal L.
 \label{eq:pg-det-kernel-splitting}
\end{equation}
任意の底変更$\alpha:T\to\mathcal V_m$について、
核・像・余核および\eqref{eq:pg-det-kernel-splitting}の形成は底変更と両立する。
特に
\begin{equation}
 \alpha^*\mathcal K\simeq\ker(\alpha^*\mathcal D),
 \qquad
 \alpha^*\mathcal L\simeq
 \ker(\alpha^*\mathcal D|_{\alpha^*\mathcal E_3^{\mathrm{act}}}).
 \label{eq:pg-det-kernel-base-change}
\end{equation}
\end{lemma}
\begin{proof}
$\mathcal D$は純中心三次部分を消す。
各点の近傍で$r_m$次小行列式を可逆にし、行と列の基本変形を行う。
全ての$(r_m+1)$次小行列式は恒等的に零なので、
行列は単位行列$I_{r_m}$と零行列からなる形となる。
従って核・像・余核は局所自由で、対応する完全列は局所的に分裂する。
この形は非被約な底への変更でも保たれる。
$\mathcal E_3^{\mathrm{act}}$の階数は$r_m+1$であるから、その核は階数一である。
純中心部分への射影$\rho$が分裂を与える。
\end{proof}

\begin{remark}[通常の中心の形成との区別]
\label{rem:pg-det-center-base-change}
この補題が扱うのは$\hbar$で\emph{割った}交換子の核である。
$h_T=\alpha^*\hbar$、$\mathcal D_T=\alpha^*\mathcal D$と書くと、
通常の中心については
\begin{equation}
 \begin{aligned}
 Z(\mathcal A_T)_3&=\ker(h_T\mathcal D_T),\\
 Z(\mathcal A_T)_3/\alpha^*\mathcal K
 &\simeq
 \operatorname{Ann}_{\OO_T}(h_T)
 \otimes_{\OO_T}\operatorname{im}\mathcal D_T.
 \end{aligned}
 \label{eq:pg-det-ordinary-center-base-change}
\end{equation}
実際、$\mathcal D_T$の像への全射は局所的に分裂しており、
中心元の像はちょうど$h_T$で消える部分である。
$h_T$が非零因子なら両方の核は一致する。
しかし$h_T=0$では代数は可換なので通常の中心は$\mathcal E_{3,T}$全体であり、
割った交換子の核は古典Casimir部分
$\mathcal C_{m,T}+\OO_T f$である。
従って\eqref{eq:pg-det-kernel-base-change}を、通常の中心が
任意の底変更と両立するという主張に読み替えてはならない。
\end{remark}

\begin{proposition}[Cramer規格化の一意性と遷移]
\label{prop:pg-det-normalization-transition}
全次数三でactive変数を含む単項式$w_i$を固定し、
その係数写像を$\varepsilon_i$、$f$の$w_i$係数を$\alpha_i$とする。
$\alpha_i$が可逆で、$\varepsilon_i|_{\mathcal L}$が自明化を与える開部分では、
Cramer中心元$\Omega_i$は次の三条件によって一意に定まる：
\begin{equation}
 \mathcal D\Omega_i=0,\qquad
 \rho(\Omega_i)=f_{\mathrm{cent}},\qquad
 \varepsilon_i(\Omega_i-f_{\mathrm{cent}})=\alpha_i.
 \label{eq:pg-det-normalization-conditions}
\end{equation}
特に、同じ$w_i$を用いる限り、可逆な小行列を作る\emph{行の選択}に依存しない。
二つの規格化が有効な共通開部分では
\begin{equation}
 \begin{aligned}
 \Omega_j&=a_{ij}\Omega_i+c_{ij},\\
 a_{ij}&=
 \frac{\alpha_j}{\varepsilon_j(\Omega_i-f_{\mathrm{cent}})}
 \in\OO^\times,\qquad
 c_{ij}=(1-a_{ij})f_{\mathrm{cent}}\in\mathcal C_m.
 \end{aligned}
 \label{eq:pg-det-affine-transition}
\end{equation}
一般に、$\mathcal K/\mathcal C_m$の基底へ写る二つの局所切断
$\Omega,\Omega'$の間には、一意な単元$a$と純中心三次式$c$による
$\Omega'=a\Omega+c$が成り立つ。
両方が法$\hbar$で$f$となるなら、$a|_{\hbar=0}=1$、$c|_{\hbar=0}=0$である。
\end{proposition}
\begin{proof}
$\Omega_i-f_{\mathrm{cent}}$は$\mathcal L$の切断であり、
同線束の自明化$\varepsilon_i$における値を指定すれば一意に決まる。
定理\ref{thm:pg-universal-determinant}のCramer式はこの条件を満たすので、
行の選択を変えても同じ元を与える。
共通開上では$\Omega_i-f_{\mathrm{cent}}$と
$\Omega_j-f_{\mathrm{cent}}$がともに$\mathcal L$の基底である。
その比を$\varepsilon_j$で計算すると\eqref{eq:pg-det-affine-transition}となる。
一般の場合も商線束での比を$a$とすれば、差は$\mathcal C_m$に属する。
法$\hbar$で$\CC f\oplus R_3$が直和であることから最後の主張を得る。
従って同じ古典極限を持つという条件だけでは、$a=1$とは結論できない。
\end{proof}

\begin{theorem}[固定した変数分解の全係数空間における中心元]
\label{thm:pg-det-global-coefficient-lift}
全三次式の係数を$\boldsymbol c$、滑らかさの判別式を$\Delta$とし、
$B_0=\CC[\boldsymbol c,\Delta^{-1}]$とおく。
従って$\mathcal T_m=\operatorname{Spec}B_0$である。
$\tau$を新しい不定元とすると、ある$Q\in B_0[\tau]$、$Q(0)=1$が存在し、
\begin{equation}
 B=B_0[\tau,Q(\tau)^{-1}],\qquad
 B_h=B_0[\hbar,Q(\hbar^2)^{-1}]
 \simeq B[\hbar]/(\hbar^2-\tau)
 \label{eq:pg-det-even-coefficient-rings}
\end{equation}
上で次が成り立つ。
固定したactive変数$x,y,z$と中心変数$u$に関する普遍代数には、
\begin{equation}
 \Omega_{\mathrm{univ}}\in Z(\mathcal A_{B_h})_3,
 \qquad
 \Omega_{\mathrm{univ}}\bmod\hbar=f_{\mathrm{univ}},
 \qquad
 \rho(\Omega_{\mathrm{univ}})=f_{\mathrm{cent}}
 \label{eq:pg-det-global-center}
\end{equation}
を満たす中心元が存在する。
そのWeyl係数は$B$に属し、$\hbar$について偶である。
さらに$\Omega_{\mathrm{univ}}-f_{\mathrm{cent}}$は商線束の大域的な基底となる。
従って全ての滑らかな$f$について、零切断全体を含む\emph{一つの}係数開集合上に、
係数が正則な中心三次元を同時に選べる。
この選択の正準性や座標変換に関する同変性は主張しない。
同変な選択の存在については、次節の
定理\ref{thm:pg-equivariant-levi-center}と
定理\ref{thm:pg-equivariant-parabolic-center}で扱う。
\end{theorem}
\begin{proof}
$\mathcal V_m$の補集合を閉集合として定めるイデアルを
$I\subset B_0[\hbar]$とする。
補集合は零切断と交わらないので$I+(\hbar)=B_0[\hbar]$であり、
$p\in I$、$p(0)=1$を選べる。
$p(\hbar)p(-\hbar)$は偶多項式であるから、
これを$Q(\hbar^2)$と書く。
そのprincipal開は$\mathcal V_m$に含まれ、零切断全体を含む。
必要なら最初に$\hbar\mapsto-\hbar$で移した開部分との共通部分を取っておく。

Weyl基底では$\mathcal D$の全成分が$\hbar$について偶であるから、
$\mathcal D$は$B$上の行列へ降りる。
$B_h=B\oplus\hbar B$は$B$上忠実平坦である。
従って$r_m$次小行列式が生成するイデアルは$B$上でも単位イデアルであり、
全ての$(r_m+1)$次小行列式は零である。
補題\ref{lem:pg-det-divided-kernel}の行列操作を$B$上で行うと、
active部分の核$L_B$は階数一の有限射影$B$-加群となる。
$B$は多項式環の局所化なのでUFDであり、
$\operatorname{Pic}(B)=0$である\cite[Tag~0BCH]{PStacks}。
従って$L_B$の基底$\ell$を選べる。

$Q(0)=1$より$B/(\tau)=B_0$である。
古典Casimir核の計算から、$L_B/(\tau)$の基底は
$f_{\mathrm{univ}}-f_{\mathrm{cent}}$である。
実際、この切断は全ての滑らかな$f$で非零なので線束の基底となる。
従ってある$b\in B_0^\times$について
\begin{equation}
 \ell\bmod\tau=b(f_{\mathrm{univ}}-f_{\mathrm{cent}}).
 \label{eq:pg-det-global-generator-normalization}
\end{equation}
$b^{-1}$を$B$の単元とみて$\ell$を$b^{-1}\ell$へ換え、
$\Omega_{\mathrm{univ}}=f_{\mathrm{cent}}+\ell$とすればよい。
$\mathcal D\ell=0$は全交換子の消滅を意味し、
$B$に属する係数は$\hbar$について偶である。
\end{proof}

\begin{example}[規格化の変更と中心三次イデアル]
\label{ex:pg-det-distinct-normalized-ideals}
$f=(x^3+y^3+z^3)/3+xyz+u^3$を取る。
これは滑らかなHesse三次と$u^3$の和なので全三次超曲面は滑らかである。
$S=x^3+y^3+z^3$、$T=\sym(xyz)$、$\tau=\hbar^2$とし、
\begin{equation}
 U=1+\frac{\tau}{2}+\frac{7\tau^2}{48},\qquad
 W=1-\frac{5\tau}{6}+\frac{7\tau^2}{48},\qquad
 U-W=\frac{4\tau}{3}
 \label{eq:pg-det-hesse-transition-coefficients}
\end{equation}
とおく。定理\ref{thm:pg-universal-hesse-mixed}の$t=1$の式により、
$US/3+WT+u^3$は中心元である。
$U,W$が可逆な原点近傍で、
\begin{equation}
 \Omega_x=u^3+\frac S3+\frac WU T,\qquad
 \Omega_T=u^3+\frac U{3W}S+T
 \label{eq:pg-det-hesse-normalized-centers}
\end{equation}
も中心元となる。
命題\ref{prop:pg-det-normalization-transition}の一意性により、
これらはそれぞれ$x^3$の係数を$1/3$、$xyz$の係数を$1$に固定したCramer規格化である。
両方の古典極限と純中心係数は同じだが、
\begin{equation}
 \Omega_T=\frac UW\Omega_x+
 \left(1-\frac UW\right)u^3
 \label{eq:pg-det-hesse-affine-transition}
\end{equation}
であり、その差は純中心項だけではない。
通常のWeyl基底も有効な共通開へ縮めると、
$\hbar\ne0$では、自然な周囲代数の同一視の下で
$(\Omega_x)$と$(\Omega_T)$は異なるイデアルである。
実際、両イデアルが一致すれば次数三部分を比較して
$\Omega_T=b\Omega_x$となる。
$u^3$係数より$b=1$だが、$S$の係数は$U/W\ne1$なので矛盾する。
これは中心商の抽象的な代数同型の不存在を主張する例ではなく、
affineな遷移だけでは同じ中心イデアルを得ないことを示す。
\end{example}

\begin{remark}[ブロックを保つ座標変換]
\label{rem:pg-det-block-covariance}
以下の比較は、双方の構成が定義される共通開部分上で行う。
定理で選んだprincipal開そのものの座標変換に関する不変性は主張しない。
$M\in\mathrm{GL}_3(\CC)$、$N\in\mathrm{GL}_m(\CC)$、
$C\in\operatorname{Mat}_{3\times m}(\CC)$、$\lambda\in\CC^\times$に対し、
\begin{equation}
 \begin{aligned}
 X&=Mx+Cu,\qquad U=Nu,\\
 f'(X,U)&=\lambda^{-1}
 f\bigl(M^{-1}(X-CN^{-1}U),N^{-1}U\bigr),\qquad
 \hbar'=\lambda\det(M)\hbar
 \end{aligned}
 \label{eq:pg-det-block-coordinate-change}
\end{equation}
とおくと、代入は$A_{f',\hbar'}\simeq A_{f,\hbar}$を与える。
中心変数は可換であり、Weyl対称化はこの線形代入と可換する。
三重ベクトルの変換則から、代入を$\Phi$と書けば
\begin{equation}
 \Phi^{\oplus3}(\mathcal D_{f',\hbar'}p')
 =\frac1{\lambda\det M}\,
 M\mathcal D_{f,\hbar}(\Phi p')
 \label{eq:pg-det-divided-matrix-covariance}
\end{equation}
となる。従って$\mathcal K$、純中心部分、およびその商線束はこの変換と両立する。
Poisson係数を$\eta=\nu\hbar$と分けて記す場合は、
$\nu'=\lambda\det(M)\nu$とすることで共通の変形変数$\hbar$を固定できる。
一方、$f'$へ規格化した中心元を$\Omega'$とすると、
$\lambda\Phi(\Omega')$と$f$へ規格化した$\Omega$の間に保証されるのは
$\lambda\Phi(\Omega')=a\Omega+c$という遷移である。
特に中心平行移動は純中心部分への射影$\rho$を保たないため、
この代数の変換則から中心元の選択の同変性までを結論することはできない。
\end{remark}

\begin{remark}[係数空間上の存在とモジュライへの降下]
\label{rem:pg-det-global-descent-scope}
定理\ref{thm:pg-det-global-coefficient-lift}は、固定した変数分解の係数空間上で
中心元を一括して選べることを示す。
この底では商線束の非自明性を障害とする必要はない。
しかし例\ref{ex:pg-det-distinct-normalized-ideals}の通り、
別々に規格化した中心商が自然な周囲代数の同一視で貼り合わさるとは限らない。

純中心部分との交わりが零で、商線束へ同型に写る中心三次直線を選ぶことは、
完全列$0\to\mathcal C_m\to\mathcal K\to\mathcal L\to0$の分裂を選ぶことと同値である。
その選択空間は$\mathcal Hom(\mathcal L,\mathcal C_m)$の作用するaffine束であり、
中心三次イデアルの選択を記述する。
モジュライへ降下するには、この選択を座標変換に関して同変にし、
零切断で$\CC f$を選ぶ条件および遷移のcocycleと両立させる必要がある。
複数ブロックでは、ブロック分解と純中心三次の配分の変化との両立も必要である。
これらは本定理のfactorialな係数環上での基底選択とは別の条件であり、
次節では変数分解を保つreductive群について、また三変数制限が滑らかな場合には
中心平行移動を含む群について、これらの条件を満たす選択を構成する。
特異な三変数制限を含む全座標群への拡張と、複数ブロックの配分との整合性を含めた
全Poissonモジュライへの大域量子化の降下は主張しない。
\end{remark}

\begin{lemma}[中心三次式の配分による全ブロックの滑らか化]
\label{lem:pg-central-splitting-smooth}
$X=V(F)\subset\PP^{N-1}$が滑らかで、
正規形定理\ref{thm:pg-normal}によって
\begin{equation}
 F=\sum_{s=1}^k\widetilde f_s(x_s,u)+c(u),
 \qquad
 \widetilde f_s(x_s,0)\in\CC[x_s]_3,\quad
 c(u)\in\CC[u]_3
 \label{eq:pg-det-central-splitting}
\end{equation}
と書けるとする。各$\widetilde f_s$は
純中心三次部分を含まないように選ぶ。
このとき中心三次部分を
\begin{equation}
 c(u)=c_1(u)+\cdots+c_k(u),\qquad
 f_s(x_s,u)=\widetilde f_s(x_s,u)+c_s(u)
 \label{eq:pg-det-center-distribution}
\end{equation}
と配分して、\emph{全ての}
$V(f_s)\subset\PP(W_s\oplus U)$が
滑らかになるようにできる。
この配分は$\partial_{x_s}f_s$を変えないため、
元のPoisson構造を一切変えない。
\end{lemma}
\begin{proof}
まず$u=0$の座標平面$\PP(W_s)$を考える。
そこでは$f_s=\widetilde f_s$、
その$x_s$偏微分は$F$の対応偏微分と一致する。
他のactive変数$x_t$の偏微分は、
$x_t=0,u=0$では斉次次数三のため零である。
また$f_{s,u}$はこの平面上で$F_u$と一致し、
$c_s$の偏微分はそこで零である。
従って$\PP(W_s)$上で$f_s$が特異なら、
同じ点で$F$も特異となり矛盾する。
ゆえに$f_s$は全ての$c_s$に対し、
$u=0$上では滑らかである。

次に$m>0$とする。
$u\ne0$の開集合では、中心三次式$\CC[u]_3$の
線形系は基点自由である。
従って$\langle\widetilde f_s,\CC[u]_3\rangle$が定める
射影線形系に標数零のBertini定理を適用すると、
一般の元はこの開集合上で滑らかである。
その非空開条件と$\widetilde f_s$の係数が非零という
非空開条件を同時に満たし、その係数を1に規格化する。
これにより、固定した$\widetilde f_s$へ一般の$c_s(u)$を足す
アフィン線形系についても、
$u\ne0$でも滑らかな$f_s$を与える$c_s$の集合は
非空Zariski開となる。
$k=1$なら$f_1=F$なので既に滑らか。
$k\ge2$では
\begin{equation}
 \mathcal C=
 \{(c_1,\ldots,c_k)\in(\CC[u]_3)^k:
 \sum_sc_s=c\}\simeq(\CC[u]_3)^{k-1}
 \label{eq:pg-det-distribution-space}
\end{equation}
は既約なアフィン空間であり、
各$c_s$への射影は全射である。
従って各$f_s$の滑らかさを課す非空開条件を
$\mathcal C$上で同時に満たせる。
$m=0$では中心三次式はなく、
$F$の滑らかさから各純active$f_s$が既に滑らかである。
最後の主張は$\partial_{x_s}c_s=0$から明らかである。
\end{proof}

\begin{example}[三変数制限が恒等的零でも中心元を得る]
\label{ex:pg-det-zero-ternary}
$m=3$、$u=(u,v,w)$とし
\begin{equation}
 f=u x^2+v y^2+w z^2+u^3+v^3+w^3
 \label{eq:pg-det-zero-ternary-example}
\end{equation}
を考える。$f(x,y,z,0,0,0)=0$である。
しかし$2ux=x^2+3u^2=0$から$x=u=0$
（$y,v$、$z,w$についても同様）なので、
$V(f)\subset\PP^5$は滑らか。
古典交換子行列$S_3\to S_4^{\oplus3}$は
\begin{equation}
 \dim S_3=56,\qquad
 \dim R_3=10,\qquad
 \operatorname{rank}\mathcal D_{f,0}=45
 \label{eq:pg-det-zero-ternary-rank}
\end{equation}
を満たす。これは補題
\ref{lem:pg-det-classical-kernel}の帰結であり、
\path{computations/verify_singular_cubic_classical_kernel.py}
でも有理数演算で厳密検算できる。
しかもこの例ではCramer式を計算するまでもなく、
\begin{equation}
 [y,z]=2\hbar ux,\quad
 [z,x]=2\hbar vy,\quad
 [x,y]=2\hbar wz,
 \qquad
 \Omega_{f,\hbar}=f
 \label{eq:pg-det-zero-ternary-explicit}
\end{equation}
が\emph{全ての$\hbar$で厳密に中心}である。
実際
$[vy^2+wz^2,x]
=-2\hbar vw(yz+zy)+2\hbar vw(zy+yz)=0$で、
他の生成元との可換性も循環置換で従う。
これは$f(x,y,z,0)=0$であっても
代数的中心元が実際に存在することの直接的な例であり、
一般の場合の存在は定理\ref{thm:pg-universal-determinant}による。
\end{example}

\begin{remark}[先行研究と確実性]
\label{rem:pg-det-scope}
Hesse型・Fermat型では
定理\ref{thm:pg-universal-hesse-mixed}と
定理\ref{thm:pg-fermat-mixed-universal}の
短い多項式閉公式があり、これらを
普遍交換子行列のgeneric rank恒等式の入力とした。
本定理は、それらの中心性を代数的恒等式として用い、
正規形の各ブロックの全体の滑らかさだけから
\emph{特異制限も含む}代数的中心元を構成する。
中心式は通常は有限行列式の\emph{有理式}であって、
特異な制限の全てに共通な短い多項式閉式が存在するとは
主張しない。全体を滑らかにする配分の選択も一般に非一意。
新規性の全面的な先行研究調査は未完了である。
\end{remark}
